\chapter{How they pick you}

Attackers do not pick randomly. They pick people who are
easy to find, easy to predict, and able to move money
alone.

\section{Individual risk}

Public signaling is the loudest amplifier. Founders,
KOLs, traders, and engineers who post fundraising
numbers, PnL, or treasury size are advertising.

Lifestyle does the rest. A repeated gym, the same café,
a recognizable car, premium travel on a fixed weekday:
cheap surveillance. You do not need a skilled team to
learn when someone is alone.

Signing large transactions from an unsecured home
collapses the digital control and the body into one room.
That is the room they want.

\section{Organizational risk}

The org version is concentrated power. If one human can
authorize a material treasury movement, coercion is
efficient.

It gets worse when key management is treated as an IT
chore. Hiring, offboarding, travel, and family exposure
never enter the conversation. The signer goes to a
conference with a hot wallet and a public badge. Nobody
owned that risk because it was not a ticket in Jira.

\section{Place}

Violent crime, weak police, or corruption raises the
base rate. Travel into aggressive or opaque enforcement
adds a second problem: ``official'' coercion at a border,
a hotel, or an office. The threat model for a signer in
that airport is not the same as the threat model at home.

\section{Early signals}

\textbf{OSINT and probing.} Follower spikes, DMs that ask
how treasury works, questions about routines. The same
unfamiliar face near home or the office more than once.

\textbf{Doxxing and threats.} Harassment that names
family. Posts that glue a real identity to a wallet or a
balance. A leaked home or office address. That is a
shift from opportunistic to targeted.

Treat those as escalation, not as internet noise. Later
chapters assume you still have a window. Sometimes you
do not. The point of watching is to use the window:
lower what one person can move, change the routine, tell
the person you already picked for that call.
